\section{Conclusion}
HNE turns changes in personal rhythms into a shared representation for future-disturbance warning. Personal deviations, context responses, and relationship changes retain the same coordinate definitions across different measurement sets. The main comparisons show gains in AUROC and average precision, with relationship changes strengthening warning on both cohorts. The source-coefficient experiment further identifies a concrete transfer benefit: learned warning rules improve low-data adaptation beyond transfer of the representation basis alone. HNE thus connects individual reference behavior, multivariate precursor signals, and learning across monitoring domains.

\section*{Limitations}
Estimating how an intervention changes disturbance risk requires interventional outcomes beyond the observational endpoint studied here. Relating wearable rhythm features to endogenous circadian phase requires independent phase measurements. These are two extensions of the present prediction task.
